\section{Proof of the main theorem}\label{gen:section}
\begin{proof}[Proof of Theorem~\ref{gen:theorem}]
Suppose counterexamples exist with unbounded order. A regular digraph has no arcs between different strongly connected components: in its acyclic condensation, cut balance forces a source component to have no outgoing arc; remove it and repeat. Every component inherits the $d$ full colour matchings and has at least $d+1$ vertices. Restricting to one preserves the density lower bound and gives a strongly connected counterexample sequence whose order tends to infinity.

Apply Lemma~\ref{gad:decomposition}. Take $d/n\to{}p>0$, and let $C$ be its limiting circulation. Its directed support has no edges between distinct strong components: a source component has zero incoming flow, hence zero outgoing flow by conservation, and induction applies. Let these components be $I_1,\ldots,I_s$, and delete zero-limit physical cells. The sets $U_t=\bigcup_{i,j\in{}I_t}(B_i\cap{}A_j),\beta_t=\sum_{i\in{}I_t}a_i$ \textup{(GC.2)}. partition all but $o(n)$ physical vertices. Every real block arc remains in its own core, so each core has good internal degrees $d-o(n)$, and every good colour has $|U_t|-o(n)$ arcs there. Also $\beta_t\ge{}p$ and $\sum_t\beta_t=1$.

The normalized restriction $C[I_t]/\beta_t$ has strongly connected support, possibly periodic. In the periodic case equalize phase sizes by deleting $o(n)$ vertices, keeping each supported cell positive. Lemma~\ref{rou:power} gives positive linear families of bounded-length, phase-compatible rainbow connectors between distinct vertices of each core. This argument uses only counting and the local gaps; it precedes every absorption or final-path conclusion. Thus the connector hypothesis of Proposition~\ref{gcp:completion} is satisfied before a prefix is chosen.

It remains to check rooted completion when $\beta_t>p$. The effective density is $p_t=p/\beta_t<1$, so the physical capacity has fixed positive linear slack, surviving all $o(n)$ deletions. If $p_t>1/2$, Theorem~\ref{pf:theorem} gives the rooted suffix. If $p_t\le1/2$, the restricted approximate incidence blocks have fixed local gaps, positive limiting cells, and strongly connected quotient support; these properties and the uniform degree and colour estimates survive $o(n)$ deletions. Theorem~\ref{gaq:completion} gives the rooted suffix, including every quotient period and the one-block case.

Proposition~\ref{gcp:completion} now supplies the exceptional-colour prefix and rooted suffix. When all core masses equal $p$, its strict-capacity argument gives a suffix of density tending to one; for a single core the original density tends to one. In every case the resulting path has $d-1$ arcs, a contradiction.
\end{proof}
